\documentclass[11pt]{article}

\usepackage[a4paper, margin=1in]{geometry}
\usepackage{amsmath, amssymb}
\usepackage{xcolor}
\usepackage{listings}
\usepackage{hyperref}
\usepackage{enumitem}

\lstset{
  basicstyle=\ttfamily\small,
  language=Python,
  keywordstyle=\color{blue!70!black},
  commentstyle=\color{green!40!black}\itshape,
  stringstyle=\color{red!60!black},
  showstringspaces=false,
  breaklines=true,
  frame=single,
  framerule=0pt,
  backgroundcolor=\color{black!4},
  xleftmargin=0.5em,
  xrightmargin=0.5em,
}

\hypersetup{
  colorlinks=true,
  linkcolor=blue!60!black,
  urlcolor=blue!60!black,
}

\title{SiGMoiD: Summary of Changes\\\large from initial version (commit \texttt{c846b30}) to current \texttt{qol} branch}
\author{}
\date{}

\begin{document}
\maketitle

\section*{TL;DR}

The starting version was a working prototype with hand-rolled gradients, a
placeholder \texttt{pyproject.toml}, no tests, and a probability function that
required reading the source to understand. The current branch turns it into a
proper PyTorch package:

\begin{itemize}[leftmargin=*, itemsep=2pt]
  \item \textbf{\texttt{Model} is now a \texttt{torch.nn.Module}}, so it composes with the rest of the PyTorch ecosystem (custom optimizers, schedulers, mixed precision, \texttt{state\_dict}, GPU movement, etc.).
  \item \textbf{Hand-rolled gradient loop replaced with autograd.} Default optimizer is \texttt{Adam(lr=nu)} (was an implicit SGD). Mathematically equivalent to the original update when \texttt{SGD} is chosen explicitly.
  \item \textbf{Bug fixes}: \texttt{Selector.fit} now respects its \texttt{gpu=} argument; integer-dtype input no longer crashes; \texttt{draw\_samples} no longer mutates global NumPy RNG state.
  \item \textbf{Packaging}: real metadata, dev extras, dropped unused \texttt{torchvision} dependency, \texttt{requires-python>=3.10} (was \texttt{>=3.12}).
  \item \textbf{Quality infrastructure}: 18 unit tests, GitHub Actions CI on Python 3.10/3.11/3.12, ruff linting, optimizer benchmark.
  \item \textbf{Method rename}: \texttt{Model.train} $\rightarrow$ \texttt{Model.fit} (so \texttt{nn.Module.train(mode=True)} keeps its standard PyTorch meaning of toggling train/eval mode).
\end{itemize}

\noindent The probability function from the paper,
\[
  p_{si} \;=\; \frac{\exp(-\sum_k \beta_{sk} E_{ki})}{1 + \exp(-\sum_k \beta_{sk} E_{ki})} \;=\; \mathrm{sgm}(-z), \quad z = \beta E,
\]
is preserved \emph{verbatim}. See \S\ref{sec:sigmoid} for why this is worth flagging.

\bigskip
\hrule
\bigskip

\section{Packaging and repository hygiene}

\subsection{\texttt{pyproject.toml}}

Filled in real description, license, classifiers, and project URLs. Added a
\texttt{[project.optional-dependencies]} section with a \texttt{dev} group
(\texttt{pytest}, \texttt{ruff}, \texttt{mypy}) and a \texttt{pandas} group for
the optional \texttt{Selector.trace\_df()} helper. Loosened the Python floor
from 3.12 to 3.10 and removed the unused \texttt{torchvision} dependency.

The distribution name is now \texttt{sigmoid-py} (the importable name remains
\texttt{sigmoid}). PyPI's \texttt{sigmoid} slot is taken; \texttt{sigmoid-py}
matches the planned PyPI name in the original README.

\subsection{Layout}

\begin{lstlisting}[language=]
src/sigmoid/
  __init__.py     # now re-exports Model, Selector, __version__
  model.py
  selector.py
                  # utils.py was empty; deleted
tests/            # NEW
  test_model.py
  test_selector.py
examples/         # NEW
  bench_optimizers.py
.github/workflows/ci.yml   # NEW
docs/changes.tex  # this file
CHANGELOG.md      # NEW
\end{lstlisting}

\subsection{Top-level imports}

\textbf{Before:}
\begin{lstlisting}
from sigmoid.model import Model
from sigmoid.selector import Selector
\end{lstlisting}

\textbf{Now:}
\begin{lstlisting}
from sigmoid import Model, Selector
\end{lstlisting}

This is a one-line change to \texttt{src/sigmoid/\_\_init\_\_.py} but it's the
form most users will type, so it matters.

\bigskip
\hrule
\bigskip

\section{\texttt{Model} ported to \texttt{nn.Module}}
\label{sec:nnmodule}

This is the biggest structural change. \texttt{Model} now inherits from
\texttt{torch.nn.Module}. To understand what that buys us, here's a tour of
the new pieces.

\subsection{What is \texttt{nn.Module}?}

\texttt{nn.Module} is PyTorch's base class for anything with \emph{learnable
parameters}. By inheriting from it, our \texttt{Model} automatically gains:

\begin{itemize}[leftmargin=*, itemsep=2pt]
  \item \texttt{model.parameters()} -- iterator over all learnable tensors. Optimizers consume this.
  \item \texttt{model.to(device)} -- moves all parameters and buffers to a device in one call.
  \item \texttt{model.state\_dict()} / \texttt{model.load\_state\_dict(...)} -- save/load weights.
  \item \texttt{model.train()} / \texttt{model.eval()} -- toggles training/evaluation mode (used by dropout, batchnorm, etc.).
  \item Compatibility with \texttt{nn.DataParallel}, \texttt{torch.compile}, \texttt{torch.amp} (mixed precision), and almost any PyTorch ecosystem package.
\end{itemize}

\subsection{\texttt{nn.Parameter} vs.\ a plain tensor}

In the original code:
\begin{lstlisting}
beta = torch.normal(...)            # plain tensor
energy = torch.normal(...)          # plain tensor
\end{lstlisting}

In the new code:
\begin{lstlisting}
self.beta = nn.Parameter(torch.empty(s, self.k).normal_(mean, std))
self.energy = nn.Parameter(torch.empty(self.k, i).normal_(mean, std))
\end{lstlisting}

A \texttt{nn.Parameter} is just a tensor with a flag that says ``I am
learnable.'' When you assign one as an attribute on an \texttt{nn.Module}, the
module silently registers it in an internal dict. That's how
\texttt{model.parameters()} knows what to return, and how \texttt{state\_dict}
knows what to save.

\subsection{\texttt{register\_buffer} for non-learnable tensors}

\begin{lstlisting}
self.register_buffer("raw", torch.from_numpy(np.asarray(data)).float())
\end{lstlisting}

\texttt{raw} (the data matrix) is not learnable, but we want it to move when
the user calls \texttt{model.to("cuda")}. Registering it as a buffer tells the
module: ``this is mine, take it with you when you move, but don't try to
optimize it.''

If we'd just done \texttt{self.raw = torch.tensor(...)}, the tensor would
\emph{stay on CPU} when the user called \texttt{model.to("cuda")}, and
training would crash with a device mismatch. Buffers fix that.

\subsection{The new training loop (autograd)}

\textbf{Before} -- hand-rolled gradient ascent:
\begin{lstlisting}
for _ in range(its):
    diff = prob_sigma - sigma
    beta_init = beta.clone()
    beta = beta + nu * (diff @ energy.T)
    energy = energy + nu * (beta_init.T @ diff)
    prob_sigma = self._sigma_probability(beta, energy)
\end{lstlisting}

\textbf{Now} -- standard autograd:
\begin{lstlisting}
for _ in range(its):
    optimizer.zero_grad()                       # 1. clear leftover gradients
    prob = self.forward()                       # 2. forward pass
    loss = F.binary_cross_entropy(              # 3. compute loss
        prob, self.raw, reduction="sum"
    )
    loss.backward()                             # 4. autograd computes
                                                #    d(loss)/d(every param)
    optimizer.step()                            # 5. optimizer applies
                                                #    the update rule
\end{lstlisting}

The four-step pattern (\texttt{zero\_grad}, \texttt{forward}, \texttt{loss.backward}, \texttt{optimizer.step}) is the canonical PyTorch training loop. Every PyTorch package and tutorial uses it. By adopting it we get:

\begin{itemize}[leftmargin=*, itemsep=2pt]
  \item Any optimizer (Adam, AdamW, RMSprop, LBFGS, $\ldots$) is a one-liner.
  \item Any LR scheduler plugs in via \texttt{scheduler.step()} after \texttt{optimizer.step()}.
  \item Mixed precision via \texttt{torch.amp.autocast} works without any model-side changes.
\end{itemize}

We proved (see \texttt{CHANGELOG.md} and the math derivation we worked
through) that \texttt{torch.optim.SGD(lr=nu)} on \texttt{loss = sum BCE}
produces parameter updates \emph{identical} to the original hand-rolled loop,
down to floating-point noise ($\sim 10^{-8}$ max diff after 500 iterations).
So nothing changed about the algorithm; we just expressed it in the standard
form.

\subsection{Why \texttt{Model.train} was renamed to \texttt{Model.fit}}

\texttt{nn.Module} already defines \texttt{train(mode=True)} as the method
that toggles train/eval mode. Downstream PyTorch packages call this. If we
overrode it to mean ``run gradient descent,'' those packages would silently
break. So our training method is now \texttt{fit()} (also more
\texttt{sklearn}-like).

\bigskip
\hrule
\bigskip

\section{The \texttt{@property} decorator}
\label{sec:property}

You asked specifically about this one, so a dedicated section.

\subsection{What it does}

\texttt{@property} turns a method into something that \emph{looks like an
attribute} from the caller's perspective. Compare:

\begin{lstlisting}
# Without @property:
class Model:
    def total_params(self):
        s, i = self.raw.shape
        return s * self.k + self.k * i

m.total_params()        # must call with parentheses
\end{lstlisting}

\begin{lstlisting}
# With @property:
class Model:
    @property
    def total_params(self):
        s, i = self.raw.shape
        return s * self.k + self.k * i

m.total_params          # NO parentheses; looks like an attribute
\end{lstlisting}

The behaviour is identical -- the body still runs every time you access it --
but the \emph{interface} is cleaner. You read \texttt{m.total\_params} the
same way you'd read \texttt{m.k}.

\subsection{Why we use it in the new \texttt{Model}}

Three places:

\begin{lstlisting}
@property
def total_params(self) -> int:
    """Total number of trainable parameters in the model."""
    s, i = self.raw.shape
    return s * self.k + self.k * i
\end{lstlisting}

In the old code this was \texttt{self.total\_params = self.param\_counter()}
set once in \texttt{\_\_init\_\_}. That works, but if the user ever changes
\texttt{self.k} (or we extend the model later), the cached value goes stale.
A \texttt{@property} computes it on demand from \texttt{self.raw.shape} and
\texttt{self.k}, so it's always correct. Cheap to recompute (two
multiplications), no caching needed.

\begin{lstlisting}
@property
def prob_estimates(self) -> torch.Tensor | None:
    if not self._fitted:
        return None
    with torch.no_grad():
        return self.forward().detach().cpu()
\end{lstlisting}

In the old code, \texttt{self.prob\_estimates} was a tensor that got assigned
once at the end of \texttt{train()}. With autograd, \texttt{beta} and
\texttt{energy} change every step, so the ``current'' probability matrix is
always derivable from them. The property says: \emph{when someone asks for
\texttt{model.prob\_estimates}, recompute it on the fly from the current
parameters}. This means if the user runs \texttt{fit()} again, the property
returns the new probabilities -- no risk of stale state.

The \texttt{with torch.no\_grad():} block tells PyTorch ``don't bother
building the autograd graph for this computation, I'm just inspecting'' --
saves memory.

Same idea for:
\begin{lstlisting}
@property
def model_params(self) -> list[torch.Tensor] | None:
    if not self._fitted:
        return None
    return [self.beta.detach().cpu(), self.energy.detach().cpu()]
\end{lstlisting}

\subsection{When to use it}

Rule of thumb: use \texttt{@property} for ``computed views'' on internal
state -- things that \emph{feel} like attributes but are derived from other
attributes. Don't use it for slow computations (the user won't expect
attribute access to take a long time).

\bigskip
\hrule
\bigskip

\section{Other Python idioms introduced}

\subsection{\texttt{from \_\_future\_\_ import annotations}}

\begin{lstlisting}
from __future__ import annotations
\end{lstlisting}

This makes Python treat all type hints as strings (deferred evaluation).
Practical benefit: you can write \texttt{def fit(...) -> Model:} inside the
\texttt{Model} class itself without Python complaining ``\texttt{Model} isn't
defined yet.'' Without this import you'd have to write \texttt{-> "Model"} in
quotes.

\subsection{Type hints with \texttt{|} syntax}

\begin{lstlisting}
def fit(self, seed: int | None = None) -> Model:
\end{lstlisting}

\texttt{int | None} means ``either an int or \texttt{None}.'' This is the
modern (Python 3.10+) replacement for the older \texttt{Optional[int]} from
the \texttt{typing} module. Strictly cosmetic; same meaning.

\subsection{\texttt{np.random.default\_rng} instead of \texttt{np.random.seed}}

\textbf{Before:}
\begin{lstlisting}
np.random.seed(seed)
indices = np.random.randint(...)
samples = np.random.binomial(...)
\end{lstlisting}

\textbf{Now:}
\begin{lstlisting}
rng = np.random.default_rng(seed)
indices = rng.integers(...)
samples = rng.binomial(...)
\end{lstlisting}

\texttt{np.random.seed} mutates a \emph{global} state shared by every
NumPy call in the program. If a user did:
\begin{lstlisting}
np.random.seed(7)
x = np.random.rand()        # uses seed 7
samples = model.draw_samples(seed=42)  # silently re-seeds to 42!
y = np.random.rand()        # NOT what user expected
\end{lstlisting}
\texttt{model.draw\_samples} would silently change the user's downstream
random numbers. Using a local generator (\texttt{rng}) keeps our randomness
isolated.

\subsection{Lazy imports}

\begin{lstlisting}
def trace_df(self):
    try:
        import pandas as pd
    except ImportError as exc:
        raise ImportError("trace_df() requires pandas. ...") from exc
    return pd.DataFrame(self.trace, ...)
\end{lstlisting}

\texttt{pandas} is a heavy dependency (loads NumPy, dateutil, lots of C
extensions). We don't want to pay that import cost on \texttt{import sigmoid}
just because of one optional helper. The lazy import means pandas is only
loaded if-and-when the user actually calls \texttt{trace\_df()}.

\subsection{\texttt{logging} instead of \texttt{print}}

\begin{lstlisting}
logger = logging.getLogger(__name__)
...
logger.info("Training candidate %d/%d ...", done, total, ...)
\end{lstlisting}

\texttt{print} is unconditional and goes to stdout. \texttt{logging} can be
filtered, redirected, formatted, or silenced by the user without modifying
our code. This is what every library is expected to do; \texttt{print}
inside a library is considered a code smell.

\bigskip
\hrule
\bigskip

\section{The probability function (a math note)}
\label{sec:sigmoid}

This was the hardest gotcha during the refactor.

The paper's per-cell probability is
\[
  p_{si} \;=\; \frac{e^{-z_{si}}}{1 + e^{-z_{si}}}, \qquad z_{si} \;=\; \sum_k \beta_{sk} E_{ki}.
\]

By a one-line algebraic manipulation,
\[
  \frac{e^{-z}}{1 + e^{-z}} \;\stackrel{\times e^z / e^z}{=}\; \frac{1}{1 + e^{z}} \;=\; \mathrm{sgm}(-z),
\]
so the paper formula is the standard sigmoid evaluated at $-z$, \emph{not}
at $+z$.

The original code used a hand-rolled \texttt{\_sigmoid\_transform(x)} that
implemented $\mathrm{sgm}(-x)$ in a numerically-stable way (one branch for
$x \ge 0$, another for $x < 0$, both avoiding overflow). The current
\texttt{forward()} method is now a one-liner:
\begin{lstlisting}
def forward(self) -> torch.Tensor:
    return torch.sigmoid(-(self.beta @ self.energy))
\end{lstlisting}
which is mathematically and numerically equivalent (\texttt{torch.sigmoid}
applies the same overflow-avoidance trick internally).

The mistake to never make: using \texttt{torch.sigmoid(beta @ energy)}
without the negation. That computes the \emph{complement} $1 - p$ and
silently swaps the roles of 0 and 1 in the data, breaking training without
any error. The unit test \texttt{test\_training\_actually\_minimizes\_nll}
exists specifically to guard against this regression -- if someone removes
the negation, that test fails.

\bigskip
\hrule
\bigskip

\section{Optimizer change: SGD $\rightarrow$ Adam by default}

The autograd port lets us swap optimizers freely. We benchmarked
several on a 200$\times$60 binary dataset with true latent dim 4
(\texttt{examples/bench\_optimizers.py}, 5 seeds, 1500 iterations):

\begin{center}
\begin{tabular}{lrrr}
\hline
optimizer & final NLL & iters to 99\% & wall (s) \\
\hline
SGD lr=$10^{-3}$ (old default) & $5178 \pm 188$ & 913 & 0.83 \\
SGD lr=$10^{-2}$               & $5126 \pm 183$ & \textbf{243} & \textbf{0.48} \\
SGD lr=$10^{-1}$               & \textbf{diverges} & --- & --- \\
Adam lr=$10^{-3}$              & $5346 \pm 168$ & 1407 & 0.64 \\
\textbf{Adam lr=$10^{-2}$ (new default)} & $5139 \pm 187$ & 511 & 0.69 \\
Adam lr=$10^{-1}$              & $\mathbf{5044 \pm 183}$ & 494 & 0.69 \\
AdamW lr=$10^{-2}$             & $5142 \pm 187$ & 498 & 0.69 \\
\hline
\end{tabular}
\end{center}

Why Adam is the new default:
\begin{itemize}[leftmargin=*, itemsep=2pt]
  \item Robust to learning-rate choice. SGD blows up at lr=$0.1$; Adam at the same lr is the \emph{best} run.
  \item Slightly lower final NLL than SGD on every learning rate we tried.
  \item Standard ``don't make me think'' choice for users who haven't tuned anything.
\end{itemize}

\noindent \textbf{Note for throughput-conscious users}: plain SGD with a
known-good \texttt{nu} reaches the same plateau in roughly half the wall-time
because there's no momentum/variance state to update each step. If you have
already tuned \texttt{nu} for your problem, pass:
\begin{lstlisting}
opt = torch.optim.SGD(model.parameters(), lr=nu)
model.fit(its=500, optimizer=opt, seed=42)
\end{lstlisting}
This is also mathematically equivalent to the original SiGMoiD update.

\bigskip
\hrule
\bigskip

\section{Bug fixes}

\begin{description}[leftmargin=*, itemsep=4pt]
  \item[\texttt{Selector.fit} ignored its \texttt{gpu} argument.] The
        argument was accepted but never forwarded to \texttt{Model.train}.
        Users passing \texttt{gpu=False} still trained on GPU. Fixed.

  \item[Integer-dtype input crashed BCE.]
        \texttt{torch.from\_numpy(data)} preserves the user's dtype. If they
        passed an \texttt{int64} array (very common for binary data),
        \texttt{F.binary\_cross\_entropy} failed with a cryptic error at
        evaluation time. Fixed by casting to float in \texttt{\_\_init\_\_}.

  \item[\texttt{draw\_samples} mutated global NumPy RNG state.]
        See \S4 above. Fixed by switching to a local generator.

  \item[\texttt{Model.train(k=...)} silently overrode \texttt{latent\_dim}.]
        The original \texttt{train()} method had a \texttt{k=1} kwarg that
        shadowed the latent dimension passed to \texttt{\_\_init\_\_}. A user
        following the docstring example would silently train a $k{=}1$ model
        regardless of what they constructed. Removed.

  \item[Docstring typos.] e.g.\ \texttt{""Compute the log-likelihood} (stray
        quote), \texttt{"""\_summary\_} (boilerplate placeholder).
\end{description}

\bigskip
\hrule
\bigskip

\section{Testing and CI}

\subsection{Tests}

There are now 18 unit tests in \texttt{tests/}, covering:

\begin{itemize}[leftmargin=*, itemsep=2pt]
  \item \textbf{API contract}: \texttt{fit} returns \texttt{self}; pre-fit calls to \texttt{log\_likelihood} or \texttt{draw\_samples} raise \texttt{ValueError}; \texttt{Model} is an \texttt{nn.Module} with two registered parameters.
  \item \textbf{Determinism}: same seed gives same result.
  \item \textbf{Regression tests} for the bugs listed above (gpu flag respected, integer input accepted, NumPy RNG not mutated).
  \item \textbf{Math sanity}: \texttt{test\_training\_actually\_minimizes\_nll} asserts NLL strictly decreases over training. This catches sign errors, an inverted forward pass, or a broken optimizer wiring.
  \item \textbf{Ecosystem interop}: a custom Adam optimizer accepts our parameters; an optimizer not wired to our parameters is rejected with a clear error; \texttt{state\_dict} round-trips cleanly.
\end{itemize}

\subsection{Continuous integration}

\texttt{.github/workflows/ci.yml} runs \texttt{ruff check} and \texttt{pytest}
on every push and PR, against Python 3.10 / 3.11 / 3.12. Uses the CPU PyTorch
wheel index so CI runs are quick and free.

\subsection{Linting}

\texttt{[tool.ruff]} in \texttt{pyproject.toml} configures
\href{https://docs.astral.sh/ruff/}{ruff} with a sensible rule set
(\texttt{E}, \texttt{F}, \texttt{I}, \texttt{UP}, \texttt{B}, \texttt{SIM}).
This catches unused imports, shadowed variables, deprecated typing forms,
common bug patterns, and code that can be simplified.

\bigskip
\hrule
\bigskip

\section{What was \emph{not} changed}

The published model and its mathematical behaviour are unchanged.
Specifically:
\begin{itemize}[leftmargin=*, itemsep=2pt]
  \item The probability function $p = \mathrm{sgm}(-\beta E)$ from the paper.
  \item The AIC formula $2k - 2\,\mathcal{LL}$.
  \item The model selection logic (per-candidate seed, AIC ranking).
  \item The sampling procedure (sample $\beta$ rows uniformly, then Bernoulli with that row's probabilities).
\end{itemize}

When using \texttt{SGD(lr=nu)} the parameter trajectory matches the original
hand-rolled loop to floating-point noise. Switching to Adam changes the
trajectory but not the model class or the objective.

\section*{Where to look first}

\begin{description}[leftmargin=*, itemsep=2pt]
  \item[\texttt{CHANGELOG.md}] -- canonical list of changes; the
        ``\textbf{Note -- the SiGMoiD probability function}'' callout at the
        top is the one you should re-read whenever editing
        \texttt{Model.forward}.
  \item[\texttt{src/sigmoid/model.py}] -- the rewritten model.
  \item[\texttt{tests/test\_model.py}] -- shows how the new API is meant to be used.
  \item[\texttt{examples/bench\_optimizers.py}] -- reproducible optimizer benchmark; rerun with your own data shape if you want.
\end{description}

\end{document}
